% country: RUS
The function $\psi$ from the set $\mathbb{N}$ of positive integers into itself is defined by the equality
\[
\psi(n) = \sum_{k=1}^{n} (k, n), \qquad n \in \mathbb{N},
\]
where $(k, n)$ denotes the greatest common divisor of $k$ and $n$.
(a) Prove that $\psi(mn) = \psi(m)\psi(n)$ for every two relatively prime $m, n \in \mathbb{N}$.
(b) Prove that for each $a \in \mathbb{N}$ the equation $\psi(x) = ax$ has a solution.
(c) Find all $a \in \mathbb{N}$ such that the equation $\psi(x) = ax$ has a unique solution.
